\section*{第 \ref{ch:g5:speed-distance-time} 章 --- 速度：距离与时间}

\begin{solution}{exo:g5:speed-distance-time:1}
距离相同 --- 用时少者胜；时间相同 --- 距离长者胜。
“我跑得更远”也不能说明问题：时间不等时，
跑得更远并不能说明速度更快。
\end{solution}

\begin{solution}{exo:g5:speed-distance-time:2}
博 --- 距离相同（同一游泳池），用时更少（$25 < 30$ 秒）：
规则一。
\end{solution}

\begin{solution}{exo:g5:speed-distance-time:3}
滑板车 --- 时间相同（一小时），距离更长（$45 > 20$
千米）：规则二。
\end{solution}

\begin{solution}{exo:g5:speed-distance-time:4}
每小时行进 \qty{200}{km}。$2$ 小时内：$200 \times 2
= \qty{400}{km}$。半小时内：$200$ 的一半，即 \qty{100}{km}。
\end{solution}

\begin{solution}{exo:g5:speed-distance-time:5}
蜗牛（每小时 $5$ 米），步行者（每小时 $5$ 千米），城市公交车
（$30$），信鸽（每小时 $80$ 千米）。
\end{solution}

\begin{solution}{exo:g5:speed-distance-time:6}
一小时包含 $30$ 段 $2$ 分钟：$30 \times 100 =
\qty{3000}{m}$ --- 即速度为每小时 $3$ 千米。
\end{solution}

\begin{solution}{exo:g5:speed-distance-time:7}
汽车：$80 \times 3 = \qty{240}{km}$。骑车人：$15 \times 3 =
\qty{45}{km}$。
\end{solution}

\begin{solution}{exo:g5:speed-distance-time:8}
速度较快的旅行者的图线更陡。$2$ 小时后，
骑车人的图线位于 \qty{30}{km} 处。
\end{solution}

\begin{solution}{exo:g5:speed-distance-time:9}
逐小时计算：$30$、$60$、$90$ --- $3$ 小时后，渡轮
恰好行驶了 \qty{90}{km}：横渡需要 $3$ 小时。
\end{solution}

\begin{solution}{exo:g5:speed-distance-time:10}
每小时玛拉缩短 \qty{20}{km}，乔恩缩短 \qty{10}{km}：间距
每小时缩小 \qty{30}{km}。从 \qty{60}{km} 起，
经 $2$ 小时减为零 --- 两人在距玛拉的城镇 \qty{40}{km} 处相遇。
\end{solution}

\begin{solution}{exo:g5:speed-distance-time:11}
兔子跑了一分钟后，位于 \qty{60}{m} 处 ---
正好到达中点 --- 随即睡着了。乌龟需要 $120 \div 2 = 60$
分钟，在时刻 $60$ 到达终点。兔子在时刻 $1 + 70 =
71$ 醒来，用一分钟跑完最后 \qty{60}{m}，在时刻
$72$ 到达终点。乌龟以 $12$ 分钟获胜 --- 稳扎稳打胜过
速度快，这次有算术为证。
\end{solution}
